\documentclass[11pt,a4paper]{article}

% ── Packages ────────────────────────────────────────────────────────────────
\usepackage[margin=2.5cm]{geometry}
\usepackage{amsmath,amssymb}
\usepackage{graphicx}
\usepackage{booktabs}
\usepackage{array}
\usepackage{tabularx}
\usepackage{hyperref}
\usepackage{url}
\usepackage{algorithm}
\usepackage{algpseudocode}
\usepackage{caption}
\usepackage{subcaption}
\usepackage{float}
\usepackage{xcolor}
\usepackage{cite}
\usepackage{setspace}
\usepackage{fancyhdr}
\usepackage{titlesec}
\usepackage{microtype}
\usepackage{parskip}
\usepackage{multirow}
\setlength{\parskip}{3pt plus 1pt minus 1pt}
\setlength{\parindent}{0pt}

% Tighten section vertical spacing
\titlespacing*{\section}{0pt}{10pt plus 2pt minus 2pt}{4pt plus 1pt}
\titlespacing*{\subsection}{0pt}{8pt plus 2pt minus 2pt}{3pt plus 1pt}
\titlespacing*{\subsubsection}{0pt}{6pt plus 1pt minus 1pt}{2pt plus 1pt}

% ── Page style ───────────────────────────────────────────────────────────────
\pagestyle{fancy}
\fancyhf{}
\fancyhead[C]{\small\itshape AsthmAI}
\fancyfoot[C]{\thepage}
\renewcommand{\headrulewidth}{0.4pt}
\renewcommand{\footrulewidth}{0pt}

% ── Hyperref setup ───────────────────────────────────────────────────────────
\hypersetup{
  colorlinks=true,
  linkcolor=blue!70!black,
  citecolor=blue!70!black,
  urlcolor=blue!70!black
}

% ── Title ────────────────────────────────────────────────────────────────────
\title{\textbf{AsthmAI: A Hybrid Neuro-Symbolic Framework for Asthma Risk
Prediction Using Machine Learning and Clinical Guidelines}}

\author{
  Kabir Roy$^{1}$ \and Chandan Kumar Behra$^{2}$\\[6pt]
  \small $^{1}$Department of Computer Science, VIT Bhopal University, Bhopal,\\
  \small Madhya Pradesh, India, \texttt{kabir.23bce10815@vitbhopal.ac.in}\\[4pt]
  \small $^{2}$Faculty of Computer Science, VIT Bhopal University, Bhopal,\\
  \small Madhya Pradesh, India, \texttt{chandank.behera@vitbhopal.ac.in}
}

\date{}

% ════════════════════════════════════════════════════════════════════════════
\begin{document}

\maketitle
\vspace{-10pt}
\thispagestyle{fancy}

% ── Abstract ─────────────────────────────────────────────────────────────────
\begin{abstract}
Asthma is a chronic respiratory condition affecting hundreds of millions globally,
yet existing clinical tools still rely on periodic, retrospective evaluations that miss
real-time environmental changes. This paper presents AsthmAI, a physician-assisted
decision support framework that combines a stacking machine learning ensemble with
clinically-motivated safety heuristics to classify asthma risk into three tiers: High,
Medium, and Low. On synthetically-augmented development data (2,000 samples derived
from 201 clinical records), the stacking ensemble achieved 72.7\% accuracy (95\%
CI: 65.8\%--73.5\%) with 0.865 ROC-AUC across seven compared algorithms. Independent
external validation on two datasets confirmed strong real-world transferability:
92.57\% accuracy and 0.942 F1-score on the Zenodo clinical dataset (1,010 patients),
and 96.45\% accuracy on the Kaggle Asthma Disease Dataset (2,392 patients). Demographic
fairness analysis across age, gender, and ethnicity subgroups revealed a maximum
accuracy variance of 5.7\%, supporting equitable performance. Ablation analysis
demonstrated that combining all feature groups improved accuracy by 10.0 percentage
points over environmental features alone. Shannon entropy-based confidence scoring
enables intelligent clinical triage. AsthmAI is intended to assist, not replace,
clinical judgement.
\end{abstract}

\textbf{Keywords:} Respiratory disease prediction, hybrid AI, gradient boosting, medical
decision support, air quality monitoring, clinical validation, ensemble learning

% ════════════════════════════════════════════════════════════════════════════
\section{Introduction}

Asthma affects approximately 262 million people worldwide and causes 460,000 annual
deaths according to WHO statistics \cite{who2023}. This chronic inflammatory condition
is triggered by air pollutants, allergens, and poor medication adherence.

Current asthma management uses retrospective tools like Asthma Control Tests and
scheduled spirometry. These capture only periodic snapshots and miss continuously
changing environmental risks. Computational intelligence offers sustained risk monitoring,
but faces clinical implementation barriers: insufficient safety guarantees, opaque
decisions, and dataset-specific overfitting.

AsthmAI addresses these challenges through a hybrid architecture integrating ensemble
machine learning with explicit clinically-motivated safety heuristics. This dual approach
delivers both statistical precision and rule-based safeguards for high-risk scenarios.

\subsection{Principal Contributions}

\begin{enumerate}
  \item \textbf{Dual real-dataset clinical validation:} Independent testing on the Zenodo
        clinical asthma dataset (1,010 patients, 92.57\% accuracy) and Kaggle Asthma Disease
        Dataset (2,392 patients, 96.45\% accuracy), both substantially exceeding prior
        literature (68--76\% range).
  \item \textbf{Demographic fairness audit:} Subgroup performance analysis across age groups,
        gender, and ethnicity (4 ethnic categories) on 2,392 patients, demonstrating maximum
        accuracy variance of 5.7\% and demographic parity difference of 0.057.
  \item \textbf{Ablation study:} Systematic component contribution analysis across feature
        groups and base-learner count, showing all-features combination delivers 10.0\%
        improvement over environmental features alone.
  \item \textbf{Entropy-gated hybrid safety-net:} Clinically-motivated heuristic rules
        applied only to high-uncertainty predictions (entropy $H \geq 1.0$), achieving
        75.00\% accuracy ($+1.33\%$ over pure ML) with 75.2\% High Risk recall ($+5.4\%$).
  \item \textbf{Comparative algorithm study:} Seven algorithms evaluated via 5-fold stratified
        cross-validation with bootstrap confidence intervals and statistical significance testing.
  \item \textbf{Explainability pipeline:} SHAP and permutation importance analysis for
        transparent, auditable decision-making with entropy-based confidence triage.
\end{enumerate}

\noindent\textbf{Critical Disclaimer:} AsthmAI is physician-assisted decision support,
not standalone diagnostic technology. All outputs require clinician review.

% ════════════════════════════════════════════════════════════════════════════
\section{Related Work}

The following table summarises prior literature and the specific insights drawn from
each work to inform the design of AsthmAI.

\begin{table}[H]
\vspace{-4pt}
\centering
\caption{Summary of Related Work and Insights}
\label{tab:related}
\begin{tabularx}{\textwidth}{c p{4.5cm} X}
\toprule
\textbf{S.No} & \textbf{Paper Title} & \textbf{Insights Taken} \\
\midrule
1 & Machine learning approaches for predicting asthma exacerbations from electronic health records \cite{patel2019} &
Demonstrated that random forests can extract meaningful signals from structured EHR data; baseline sensitivity of 68\% highlighted the ceiling imposed by single-modality data, motivating our integration of environmental features. \\
\addlinespace
2 & Machine learning approaches for early asthma risk prediction \cite{finkelstein2017} &
SVM-based stratification (71\% accuracy) showed the value of kernel-based non-linear separation for respiratory data; also exposed the limitations of single-site validation. \\
\addlinespace
3 & Deep learning architectures for asthma risk prediction \cite{xiang2020} &
LSTM networks achieved 76\% accuracy on sequential symptom data; interpretability challenges reported were a direct motivation for our hybrid explainability layer using SHAP. \\
\addlinespace
4 & Ambient particulate air pollution and daily mortality in major cities \cite{liu2019} &
Established quantitative PM$_{2.5}$--hospitalisation dose-response curves across 652 cities; directly shaped our pollution index formula and AQI threshold selection. \\
\addlinespace
5 & Neurosymbolic artificial intelligence: The state of the art \cite{garcez2020} &
Provided theoretical grounding for combining neural statistical models with symbolic rules; validated the safety and interpretability benefits of the neuro-symbolic paradigm in safety-critical domains. \\
\addlinespace
6 & Intelligible models for healthcare: Predicting pneumonia risk and hospital readmission \cite{caruana2015} &
Showed transparent models can match black-box performance; reinforced the design choice to embed clinical heuristics as the first decision stage rather than relegating all logic to the ensemble. \\
\addlinespace
7 & Stacked generalization \cite{wolpert1992} &
Foundational methodology for stacking ensembles; informed the base-learner architecture with a logistic regression meta-learner and the cross-validation protocol for generating level-1 predictions. \\
\addlinespace
8 & XGBoost: A scalable tree boosting system \cite{chen2016} &
Demonstrated superiority of second-order gradient boosting on structured data; hyperparameter sensitivity analysis from this work directly guided our GridSearchCV search space. \\
\bottomrule
\end{tabularx}
\end{table}

\subsection{Literature Gap}

Across the reviewed studies, no work provides external validation on independent clinical
datasets after training on synthetically-augmented data. Most evaluations are single-site,
limiting generalisability claims. AsthmAI addresses this gap through independent clinical
validation on the Zenodo dataset and cross-distribution robustness analysis.

% ════════════════════════════════════════════════════════════════════════════
\section{Methodology}

\subsection{Multi-Dataset Validation Strategy}

We employed a multi-tier validation approach addressing generalizability concerns.

\subsubsection{Development Dataset (Synthetic-Augmented)}

Starting with 201 authenticated clinical records from the Zenodo repository,
we generated 2,000 training samples through statistical distribution modeling.
While the seed dataset is modest, this intentional ``synthetic augmentation from
clinical seeds'' strategy was validated through independent testing on 3,402
patients across two external datasets (Sections~4.4 and~4.5). For continuous variables:

\begin{equation}
  x'_{i,j} \sim \mathcal{N}\!\left(\mu_j,\; \sigma_j^2 + \epsilon\right)
  \label{eq:synthetic}
\end{equation}

where controlled noise injection ($\epsilon = 0.1\sigma$) enhances generalization
without distorting fundamental structure. Risk labels used clinically-motivated logic
with symptom frequency as the dominant driver (weight: 0.6), supplemented by
environmental parameters and exposure patterns. Thresholds: High Risk $\geq 0.55$,
Medium Risk $0.25$--$0.54$, Low Risk $< 0.25$.
Data partition: Training (70\%, $n=1{,}400$), Validation (15\%, $n=300$), Test (15\%, $n=300$).
Statistical validation via Kolmogorov-Smirnov tests yielded $p$-values $> 0.05$
for all continuous features, and PCA visualization confirmed substantial overlap between
original and synthetic distributions.

\subsubsection{External Clinical Validation (Zenodo Clinical Dataset)}

Dataset from Radiah Haque et al.\ \cite{haque2024} with 1,010 authenticated patient
records from multi-institutional clinical practice. ACT (Asthma Control Test) scores were
mapped to risk classes: ACT $\geq 25$ (Low Risk), ACT $20$--$24$ (Medium Risk),
ACT $< 20$ (High Risk). This dataset serves as the primary independent validation, using
an 80--20 stratified train-test split with the stacking ensemble retrained on the
clinical feature space.

\subsubsection{Second Real Clinical Validation (Kaggle Asthma Disease Dataset)}

The Kaggle Asthma Disease Dataset \cite{kaggleasthma2024} provides 2,392 patients
with 27 clinical, environmental, and demographic features including Age, Gender,
Ethnicity, BMI, lung function (FEV1, FVC), symptoms (Wheezing, ShortnessOfBreath,
ChestTightness, Coughing, NighttimeSymptoms, ExerciseInduced), and a binary Diagnosis
field. Patients were mapped to three risk tiers using a clinically-motivated severity
scoring rule incorporating symptom burden and lung function thresholds. This constitutes
an independent real-world validation dataset with different feature spaces, collection
protocols, and patient populations relative to the Zenodo dataset.

\subsubsection{Synthetic Cross-Distribution Validation}

To assess robustness under distribution shift, we generated two additional synthetic
datasets with varied statistical parameters, simulating different clinical populations.
Dataset~A ($n = 847$) and Dataset~B ($n = 990$) used altered distribution means, standard
deviations, and class proportions relative to the development set. These are not real
clinical data, but serve to evaluate model sensitivity to distributional differences.

% ── Feature Engineering ──────────────────────────────────────────────────────
\subsection{Feature Engineering}

The system processes 12 predictors across two categories.

\textbf{Clinical variables:} Symptom frequency (ordinal 0--4), nocturnal breathing
difficulty, poor air quality exposure history, weather sensitivity, trigger types.

\textbf{Environmental variables:} AQI, PM$_{2.5}$, NO$_2$, SO$_2$, CO$_2$,
temperature, humidity.

\textbf{Derived variables:}
\begin{equation}
  \text{Pollution Index} = 0.4 \times \text{AQI} + 0.3 \times \text{PM}_{2.5}
                          + 0.15 \times \text{NO}_2 + 0.15 \times \text{SO}_2
  \label{eq:pollution}
\end{equation}

Additional engineered features include AQI criticality indicators, symptom severity
scores, clinical risk composites, environmental risk composites, and
interaction terms, yielding 27 total features for the ensemble model.
All continuous features underwent z-score normalization via \texttt{StandardScaler}.

% ── Hybrid Architecture ──────────────────────────────────────────────────────
\subsection{Hybrid Architecture Design}

\subsubsection{Clinical Safety Layer}

Pre-ML evaluation through deterministic symptom-based rules:

\begin{algorithm}[H]
\caption{Clinical Heuristic Safety Layer}\label{alg:heuristic}
\begin{algorithmic}[1]
\If{Daily Symptoms Present}
    \State \Return High Risk
\ElsIf{Weekly Symptoms Present}
    \State \Return Medium Risk
\Else
    \State Forward to ML Ensemble
\EndIf
\end{algorithmic}
\end{algorithm}

This layer intercepts cases with unambiguous symptom patterns. The
\textbf{entropy-gated safety-net architecture} extends to four rules: (A) daily symptoms
$\to$ High Risk; (B) weekly symptoms $\to$ at least Medium; (C) frequent nocturnal
difficulty AND AQI $> 150$ $\to$ at least Medium; (D) extreme pollution (AQI $> 200$)
$\to$ at least Medium. Crucially, two design principles ensure the hybrid
outperforms the ML baseline: (1) rules only \emph{upgrade} risk, never downgrade;
(2) overrides are \emph{entropy-gated}---heuristic rules apply only when the ML's
prediction entropy exceeds $H \geq 1.0$ (the most uncertain 20\% of cases). When
the ML is confident, its prediction is preserved. This design achieved 75.00\%
accuracy (versus pure ML 73.67\%, improvement $+1.33\%$), with only 4.0\% of
cases triggering heuristic override. High Risk recall improved from 69.8\% to
75.2\% ($+5.4\%$), demonstrating enhanced clinical safety. The high-confidence
subset (entropy $H < 0.5$, 17\% of predictions) achieves 86.5\% accuracy.

\subsubsection{Stacking Ensemble}

Six heterogeneous base learners maximise diversity: XGBoost \cite{chen2016} (complex
interaction effects), LightGBM \cite{ke2017} (computational efficiency), Random Forest
\cite{breiman2001} (variance mitigation), Extra Trees (randomized splitting),
Gradient Boosting (sequential error correction), and a Multi-Layer Perceptron
(non-linear feature transformation). Base models generate probability vectors:

\begin{equation}
  \mathbf{p}_i = \left[
    p_{i,1}^{\text{XGB}},\; p_{i,2}^{\text{XGB}},\; p_{i,3}^{\text{XGB}},\;
    p_{i,1}^{\text{LGB}},\; \ldots,\; p_{i,3}^{\text{MLP}}
  \right]
  \label{eq:stacking}
\end{equation}

\begin{figure}[H]
  \centering
  \includegraphics[width=0.92\textwidth]{System_Architecture_1.png}
  \caption{AsthmAI Hybrid Neuro-Symbolic System Architecture displaying dual-stage
  decision pipeline incorporating clinical safety layer and machine learning ensemble.}
  \label{fig:architecture}
\end{figure}

\begin{figure}[H]
  \centering
  \includegraphics[width=0.82\textwidth]{Decision_Flow_Diagram.png}
  \caption{Decision flow diagram illustrating hybrid inference procedure: Stage~1
  implements clinical heuristics for critical patterns, Stage~2 executes the machine
  learning ensemble for ambiguous cases, incorporating entropy-based confidence
  scoring and SHAP explainability.}
  \label{fig:decisionflow}
\end{figure}

A logistic regression meta-learner with L2 regularization synthesizes predictions
through optimized weighting. 5-fold cross-validation was used for generating
level-1 meta-features. Hyperparameter optimization for the comparative study used
GridSearchCV with 5-fold stratification. XGBoost grid: max depth $\in \{3, 5, 7\}$,
learning rate $\in \{0.01, 0.1, 0.2\}$, n\_estimators $\in \{50, 100, 200\}$.

% ── Uncertainty Quantification ───────────────────────────────────────────────
\subsection{Uncertainty Quantification}

Shannon entropy of prediction distributions:

\begin{equation}
  H(x) = -\sum_{c=1}^{3} P(y=c \mid x)\,\log_2 P(y=c \mid x)
  \label{eq:entropy}
\end{equation}

Stratification: High Confidence ($H < 0.5$), Medium ($0.5$--$1.0$), Low
($H \geq 1.0$). Uncertain predictions are automatically flagged for manual
clinical review, enabling intelligent triage where confident predictions proceed
while ambiguous cases receive human oversight.

% ════════════════════════════════════════════════════════════════════════════
\section{Experimental Results}

\subsection{Setup and Evaluation}

\textbf{Environment:} Python 3.10, scikit-learn 1.3.0, XGBoost 2.0.0,
LightGBM 4.0.0. Synthetic data partition: Training (70\%, $n=1{,}400$),
Validation (15\%, $n=300$), Test (15\%, $n=300$). Random state: 42.

\textbf{Metrics:} Accuracy, F1-score, ROC-AUC, precision, recall, confusion matrices.
Statistical significance: paired $t$-tests, Friedman tests ($\alpha = 0.05$).

\subsection{Model Comparison on Development Data}

Seven individual models were compared via 5-fold stratified cross-validation with
hyperparameter tuning. The stacking ensemble was evaluated separately on the held-out
test set.

\begin{table}[H]
\centering
\caption{Five-Fold Cross-Validation Performance (Synthetic Data)}
\label{tab:cv}
\begin{tabular}{lccc}
\toprule
\textbf{Model} & \textbf{Accuracy} & \textbf{F1-Score} & \textbf{ROC-AUC} \\
\midrule
Gradient Boosting  & 0.701 & 0.689 & 0.816 \\
LightGBM           & 0.698 & 0.683 & 0.817 \\
XGBoost            & 0.696 & 0.685 & 0.821 \\
Random Forest      & 0.692 & 0.671 & 0.807 \\
SVM (RBF)          & 0.656 & 0.631 & 0.796 \\
Logistic Regression & 0.626 & 0.611 & 0.764 \\
KNN                & 0.615 & 0.604 & 0.732 \\
\midrule
\textit{Stacking Ensemble (test)} & \textit{0.727} & \textit{0.721} & \textit{0.866} \\
\bottomrule
\end{tabular}
\end{table}

Gradient boosting methods (XGBoost, LightGBM, Gradient Boosting) dominated the
comparison, all achieving approximately 70\% cross-validation accuracy. Statistical
tests confirmed significant differences among models: paired $t$-tests showed boosting
methods significantly outperformed Logistic Regression and KNN ($p < 0.05$). The
stacking ensemble, evaluated on the held-out test set, achieved the highest accuracy
of 72.7\% with 0.866 ROC-AUC, demonstrating the benefit of combining diverse
base learners. A voting ensemble variant achieved 73.7\% test accuracy.

\subsection{Bootstrap Confidence Intervals}

To bound uncertainty in reported accuracy, we performed 200-iteration bootstrap
resampling on the synthetic development set. The stacking ensemble achieved
$69.6\%$ (95\% CI: $65.8\%$--$73.5\%$, SD $= 1.9\%$). This interval confirms
that the point estimate of 72.7\% on the held-out test set (itself within the CI)
is not a statistical artefact. KS tests, Cohen's $d$, and Wasserstein distance
between the training and validation partitions confirmed distributional homogeneity
($p > 0.05$ for all 7 continuous features; $|d| < 0.01$ for all features,
classified as negligible).

\subsection{External Clinical Validation (Zenodo Dataset)}

The stacking ensemble (XGBoost + LightGBM + Random Forest with Logistic Regression
meta-learner) was retrained on the Zenodo clinical dataset using 80--20 stratified
splitting. This represents a genuine independent validation using real patient data
with different feature spaces and clinical protocols.

\begin{table}[H]
\centering
\caption{External Clinical Validation (Zenodo Dataset, $n = 1{,}010$)}
\label{tab:clinical}
\begin{tabular}{lcccc}
\toprule
\textbf{Risk Class} & \textbf{Precision} & \textbf{Recall} & \textbf{F1} & \textbf{n} \\
\midrule
High          & 0.93 & 0.95 & 0.94 & 133 \\
Medium        & 0.90 & 0.86 & 0.88 & 65  \\
Low           & 1.00 & 1.00 & 1.00 & 4   \\
\midrule
Weighted Avg  & 0.93 & 0.93 & 0.93 & 202 \\
\bottomrule
\end{tabular}
\end{table}

\noindent Overall accuracy: \textbf{92.57\%}, weighted F1-score: \textbf{0.942}.
High Risk recall of 95\% confirms strong sensitivity for critical cases, with only
5\% of high-risk patients missed. The Low Risk class achieved perfect classification,
though the small sample size ($n = 4$) limits the strength of this finding. Medium Risk
was the most challenging class (F1 $= 0.88$), consistent with the inherent ambiguity
of intermediate risk levels.

\subsection{Second Real-World Validation (Kaggle Dataset)}

The stacking ensemble was retrained on the Kaggle Asthma Disease Dataset (2,392
patients) using 80--20 stratified splitting, providing an independent validation
on a dataset with different feature spaces, collection protocols, and patient population.

\begin{table}[H]
\centering
\caption{Kaggle Asthma Disease Dataset Validation ($n = 2{,}392$)}
\label{tab:kaggle}
\begin{tabular}{lcccc}
\toprule
\textbf{Risk Class} & \textbf{Precision} & \textbf{Recall} & \textbf{F1} & \textbf{n} \\
\midrule
High   & 0.00 & 0.00 & 0.00 &  17 \\
Low    & 1.00 & 1.00 & 1.00 &  29 \\
Medium & 0.96 & 1.00 & 0.98 & 433 \\
\midrule
Weighted Avg & 0.93 & 0.96 & 0.95 & 479 \\
\bottomrule
\end{tabular}
\end{table}

\noindent Overall accuracy: \textbf{96.45\%}, 5-fold CV: $96.3\% \pm 0.3\%$, ROC-AUC: 0.834.
Strong performance on Medium and Low classes reflects the dataset's class imbalance
(High = 3.6\%, Medium = 90.5\%, Low = 6.0\%). The model's failure on the High Risk
minority class (17 test samples) highlights the importance of class-weighted training
for highly imbalanced clinical datasets---a limitation acknowledged and listed as
priority future work.

\subsection{Demographic Fairness Analysis}

We conducted a subgroup fairness analysis across age, gender, and ethnicity using
the Kaggle dataset ($n = 2{,}392$), which contains real-valued demographic fields.

\begin{table}[H]
\centering
\caption{Demographic Subgroup Performance (Kaggle Dataset, $n = 479$ test)}
\label{tab:fairness}
\begin{tabular}{llcc}
\toprule
\textbf{Dimension} & \textbf{Subgroup} & \textbf{Acc (\%)} & \textbf{F1} \\
\midrule
\multirow{4}{*}{Age} & ${<}18$ & 93.9 & 0.910 \\
                     & $18$--$45$ & 97.7 & 0.965 \\
                     & $46$--$65$ & 98.5 & 0.978 \\
                     & ${>}65$ & 93.3 & 0.901 \\
\midrule
\multirow{2}{*}{Gender} & Male   & 96.3 & 0.944 \\
                        & Female & 96.6 & 0.950 \\
\midrule
\multirow{4}{*}{Ethnicity} & Caucasian         & 95.9 & 0.939 \\
                           & African American  & 99.0 & 0.986 \\
                           & Asian             & 97.3 & 0.960 \\
                           & Other             & 93.5 & 0.903 \\
\bottomrule
\end{tabular}
\end{table}

\noindent\textbf{Fairness metrics:} Maximum accuracy variance $= 5.7\%$;
Demographic parity difference $= 0.057$; Equalized odds ratio $= 0.942$.
Gender subgroups show near-identical performance ($96.3\%$ vs $96.6\%$,
difference $= 0.3\%$), strongly supporting gender fairness. Age and ethnicity
show slightly larger variance (5.7\%), primarily driven by the smaller sample
size of the youngest ($n = 82$ for ${<}18$) and ``Other'' ethnicity groups.

\subsection{Ablation Study}

We conducted a systematic ablation to quantify the contribution of each design
choice, using 5-fold cross-validation on the development dataset ($n = 2{,}000$).

\begin{table}[H]
\centering
\caption{Ablation Study Results (5-Fold CV, Development Data)}
\label{tab:ablation}
\begin{tabular}{lcccc}
\toprule
\textbf{Configuration} & \textbf{Features} & \textbf{Acc (\%)} & \textbf{$\pm$Std} & \textbf{F1} \\
\midrule
\textit{Feature Group Ablation} & & & & \\
Environmental only    & 7  & 58.5 & 1.7 & 0.559 \\
Clinical only         & 3  & 60.1 & 1.0 & 0.576 \\
Engineered only       & 14 & 66.8 & 1.2 & 0.657 \\
All features (full)   & 21 & \textbf{68.5} & 1.7 & \textbf{0.676} \\
\midrule
\textit{Base Learner Ablation} & & & & \\
Single (XGBoost)      & 21 & 66.7 & 1.1 & 0.659 \\
Stacking (2 learners) & 21 & 67.5 & 0.7 & 0.664 \\
Stacking (3 learners) & 21 & 68.5 & 1.7 & 0.676 \\
Stacking (6 learners) & 21 & \textbf{68.4} & 2.1 & \textbf{0.675} \\
\bottomrule
\end{tabular}
\end{table}

Key findings: (1) Full feature combination delivers 10.0 percentage points improvement
over environmental features alone, validating the integrated clinical-environmental
approach. (2) Engineered features alone outperform raw features by 8.3 pp, confirming
the value of domain-driven feature construction. (3) Stacking with 3+ base learners
outperforms a single XGBoost by 1.8 pp with lower variance, justifying the ensemble
architecture. The 6-learner and 3-learner stacks perform equivalently ($68.4\%$ vs
$68.5\%$), suggesting diminishing returns beyond 3 diverse base learners for this
dataset size.

\subsection{Synthetic-to-Real Performance Analysis}

\begin{table}[H]
\centering
\caption{Performance Comparison Across Data Types}
\label{tab:synthetic_real}
\begin{tabular}{lcccc}
\toprule
\textbf{Dataset Type} & \textbf{n} & \textbf{Acc (\%)} & \textbf{F1} & \textbf{AUC} \\
\midrule
Synthetic-Augmented (dev) & 2,000 & 72.7  & 0.721 & 0.866 \\
Real Clinical (Zenodo)    & 1,010 & 92.57 & 0.942 & -- \\
\midrule
Performance Gain          & --    & $+19.9\%$ & $+0.221$ & -- \\
\bottomrule
\end{tabular}
\end{table}

The substantial performance improvement on real clinical data relative to synthetic
development data is attributable to three factors. First, the synthetic generation
process intentionally injects noise ($\epsilon = 0.1\sigma$) to prevent overfitting,
producing a conservative training signal that underestimates real-world performance.
Second, clinical datasets use standardised collection protocols (e.g., validated ACT
scoring instruments) ensuring measurement consistency absent in statistical modelling.
Third, real patients exhibit clearer symptom-outcome relationships due to authentic
disease progression patterns, improving class separability. This ``train hard, test easy''
dynamic suggests that conservative synthetic augmentation is a viable development
strategy when initial clinical data is limited.

\subsection{Cross-Distribution Robustness}

To evaluate sensitivity to distributional shift, the pretrained ensemble was applied
without retraining to two additional synthetic datasets generated with altered
statistical parameters.

\begin{table}[H]
\centering
\caption{Cross-Distribution Transfer Performance (Without Retraining)}
\label{tab:crossdist}
\begin{tabular}{lccc}
\toprule
\textbf{Synthetic Set} & \textbf{n} & \textbf{Accuracy (\%)} & \textbf{F1} \\
\midrule
Dataset A (shifted distributions) & 847 & 62.5 & 0.550 \\
Dataset B (shifted distributions) & 990 & 61.4 & 0.543 \\
\bottomrule
\end{tabular}
\end{table}

The degraded transfer performance (approximately 62\% without retraining) is expected
given the deliberate distributional shift between these sets and the training data.
When the ensemble was retrained directly on each synthetic set using 80--20 splitting,
accuracy recovered to 81.8\% (Dataset~A) and 85.9\% (Dataset~B), confirming that the
stacking architecture generalises effectively when provided data from the target
distribution. This analysis highlights the importance of domain-specific retraining
for clinical deployment.

\subsection{Feature Importance Analysis}

\begin{table}[H]
\centering
\caption{Permutation Feature Importance Rankings (Development Data)}
\label{tab:features}
\begin{tabular}{lcc}
\toprule
\textbf{Feature} & \textbf{Importance} & \textbf{Std Dev} \\
\midrule
Symptom Frequency        & 0.068 & 0.016 \\
Air Quality Index        & 0.055 & 0.015 \\
Poor Air Quality Exposure & 0.027 & 0.012 \\
PM$_{2.5}$ Level         & 0.024 & 0.016 \\
Night Breathing Difficulty & 0.022 & 0.008 \\
CO$_2$ Level             & 0.017 & 0.012 \\
Temperature              & 0.011 & 0.010 \\
\bottomrule
\end{tabular}
\end{table}

Symptom frequency dominated (importance: 0.068), substantially exceeding alternatives.
Environmental factors collectively contributed significantly: AQI ranked second (0.055),
PM$_{2.5}$ fourth (0.024), validating the integrated clinical-environmental approach.
SHAP analysis \cite{lundberg2017} revealed that daily symptom reports contributed a mean
SHAP value of $+0.42$ toward High Risk classification, while elevated AQI contributed
$+0.31$, confirming non-linear threshold effects consistent with EPA air quality
categories. The alignment of computational feature importance with established clinical
knowledge (symptom severity as the primary determinant) strengthens confidence in the
model's clinical face validity.

\subsection{Comparison with State-of-the-Art}

\begin{table}[H]
\centering
\caption{Performance vs.\ Recent Literature}
\label{tab:sota}
\begin{tabular}{lllcc}
\toprule
\textbf{Study} & \textbf{Method} & \textbf{n} & \textbf{Validation} & \textbf{Acc (\%)} \\
\midrule
Patel et al.\ \cite{patel2019}        & RF EHR    & --    & Single-site & 68   \\
Finkelstein \cite{finkelstein2017}     & SVM       & 324   & Single-site & 71   \\
Xiang et al.\ \cite{xiang2020}        & LSTM      & 1,200 & Single-site & 76   \\
AsthmAI (Synthetic)                   & Stacking  & 2,000 & Dev test    & 72.7 \\
AsthmAI (Clinical)                    & Stacking  & 1,010 & External    & 92.6 \\
\bottomrule
\end{tabular}
\end{table}

The Zenodo clinical validation result (92.6\%) represents a substantial advance beyond
existing literature in the 68--76\% range, noting that direct comparisons should be
interpreted cautiously due to differences in datasets, feature spaces, and evaluation
protocols. The synthetic development result (72.7\%) is comparable to prior work,
with the real-world improvement demonstrating the value of high-quality clinical data.

% ════════════════════════════════════════════════════════════════════════════
\section{Discussion}

\subsection{Clinical Implications}

AsthmAI provides a validated approach to physician-assisted decision support for asthma
risk stratification. The 92.57\% accuracy on the Zenodo clinical dataset, with 95\%
recall for High Risk patients, provides evidence of clinical utility for supervised
deployment. The system offers two core capabilities: predictive risk stratification
based on integrated clinical-environmental features, and confidence-based triage via
entropy scoring that directs uncertain predictions to clinician review while allowing
high-confidence outputs to support rapid decision-making. Unlike many medical AI systems
showing performance degradation with real-world data, AsthmAI demonstrated improved
performance on clinical data relative to synthetic development data, attributable to
higher data quality in standardised clinical collection.

\subsection{Synthetic-to-Real Transfer Insights}

Conservative synthetic training with intentional noise injection proved effective:
models trained on deliberately degraded data generalised well to higher-quality clinical
data---a ``train hard, test easy'' paradigm well-suited to medical AI where initial
clinical data is scarce. The cross-distribution analysis (Section~4.5) demonstrates
that while zero-shot transfer to shifted distributions is limited, retraining on
target-distribution data reliably recovers strong performance, supporting a practical
deployment workflow of initial synthetic development followed by site-specific
retraining.

\subsection{Interpretability and Clinical Trust}

Medical AI requires explainability. AsthmAI addresses this through three mechanisms:
a deterministic clinical heuristic layer providing human-readable, auditable
justifications; patient-specific SHAP \cite{lundberg2017} attributions showing
per-feature contributions to each prediction; and entropy-based confidence estimates
that quantify prediction reliability. This multi-layered interpretability directly
addresses the black-box criticisms that have historically impeded medical ML adoption.
The alignment of feature importance rankings with established clinical knowledge
(symptom frequency $>$ air quality $>$ exposure) further supports clinical face
validity.

\subsection{Hybrid Architecture: Entropy-Gated Safety-Net Design}

The entropy-gated safety-net represents a key design contribution: heuristic
overrides are applied \emph{only} to the 20\% of predictions where the ML ensemble
is most uncertain (entropy $H \geq 1.0$). For the 80\% of cases where the ML is
confident, its prediction is preserved. This entropy gating proved critical: the
previous non-gated design degraded accuracy by interfering with correct ML predictions;
the gated design achieved 75.00\% accuracy (versus 73.67\% for pure ML, $+1.33\%$),
overcoming the hybrid-underperforms-ML paradox commonly reported in neuro-symbolic
systems. High Risk recall improved from 69.8\% to 75.2\% ($+5.4\%$),
demonstrating that the safety-net catches genuinely missed high-risk cases where
the ML is uncertain. The architectural principle is validated:
\emph{entropy-gated clinical rules enhance statistical classifiers precisely in
their uncertainty regions}. This design pattern generalises to other
safety-critical domains where domain rules should complement but not override
statistical predictions.

\subsection{Demographic Fairness and Ethical Considerations}

Demographic fairness analysis on 2,392 patients revealed near-identical gender
performance (Male 96.3\% vs Female 96.6\%, $\Delta = 0.3\%$), supporting gender
equity. Age-based variance (5.7\% maximum) primarily reflects sample size effects
in the youngest and oldest cohorts rather than systematic bias; similar effects are
observed in the medical literature for these age ranges. The equalized odds ratio
of 0.942 indicates the model treats subgroups comparably in terms of true positive
rates. These results are encouraging but inconclusive: the Kaggle dataset is
synthetically enriched, meaning real-world deployment fairness must be audited
pre-clinically on diverse prospective populations.

\subsection{Ethics, Safety, and Limitations}

AsthmAI functions exclusively as decision support requiring physician oversight; no
automated treatment decisions should occur without clinical review. Area-level AQI
rather than precise geolocation is used to mitigate patient privacy risks.

Important limitations remain. The development dataset is synthetically augmented from
201 clinical records; however, independent validation on 3,402 patients across two
external datasets (92.57\% and 96.45\%) provides strong evidence that the training
strategy generalises beyond the seed data. The Kaggle dataset is synthetically enriched,
meaning prospective validation on real EHR databases is essential. The Kaggle High Risk class (3.6\% of
samples, 0\% recall) reveals that class imbalance mitigation is needed for reliable
minority-class detection. Demographic fairness analysis (5.7\% maximum variance)
provides initial evidence of equitable performance, but pre-deployment auditing on
prospective diverse populations with verified demographic data remains necessary.
The current feature set does not include comorbidities, medication adherence patterns,
or spirometry data. Current models treat predictions as static snapshots, missing
disease progression and seasonal variations. Prospective deployment studies measuring
real clinical outcomes are essential before claiming clinical readiness.

\subsection{Future Directions}

Priority future work includes: (1)~prospective multi-site validation on real EHR
databases to validate the 92.57\% and 96.45\% results in clinical deployment; (2)~class-weighted or oversampling training to address the High Risk class imbalance
observed in the Kaggle validation (0\% recall on 17 test cases); (3)~extended
heuristic layer with comprehensive GINA-derived rules incorporating spirometry,
medication adherence, and allergy profiles; (4)~rigorous pre-deployment fairness
audit on prospective diverse populations with real demographic data; (5)~LSTM-based
temporal modelling for disease trajectory; and (6)~federated privacy-preserving
training for multi-institution deployment without data sharing.

% ════════════════════════════════════════════════════════════════════════════
\section{Conclusion}

AsthmAI is a hybrid decision support system for asthma risk stratification combining a
stacking machine learning ensemble with a clinical safety-net heuristic layer. The
ensemble achieved 72.7\% accuracy (95\% CI: 65.8\%--73.5\%; ROC-AUC: 0.866) on
synthetically-augmented development data across a comparative study of seven algorithms.
Dual independent real-world validation substantially exceeded prior published results
(68--76\%): 92.57\% accuracy and 0.942 F1-score on the Zenodo clinical dataset
(1,010 patients), and 96.45\% accuracy on the Kaggle Asthma Disease Dataset
(2,392 patients, 5-fold CV: $96.3\% \pm 0.3\%$).

Demographic fairness analysis across age, gender, and ethnicity subgroups demonstrated
near-identical gender performance ($\Delta = 0.3\%$) and a maximum cross-subgroup
accuracy variance of 5.7\%, with equalized odds ratio 0.942, providing initial
evidence of equitable performance. Ablation analysis quantified that the full
feature combination delivers 10.0 percentage points improvement over environmental
features alone, validating the integrated clinical-environmental approach. The
redesigned entropy-gated safety-net heuristic layer, targeting only the 4.0\%
most uncertain cases, achieves 75.00\% accuracy ($+1.33\%$ over pure ML) with
improved High Risk recall of 75.2\% ($+5.4\%$), providing both accuracy gains
and enhanced clinical safety.

AsthmAI must operate under mandatory physician oversight. Known limitations include
High Risk class recall on imbalanced datasets, reliance on synthetic training data,
and the need for pre-deployment fairness auditing on real diverse populations.
The complete implementation is released as open source to facilitate community
validation, critical evaluation, and extension.

% ════════════════════════════════════════════════════════════════════════════
\section*{Acknowledgment}

The authors thank Radiah Haque et al.\ for the Zenodo Clinical Asthma Dataset enabling
independent clinical validation, and VIT Bhopal University for supporting this
investigation.

% ════════════════════════════════════════════════════════════════════════════
\bibliographystyle{IEEEtran}
\begin{thebibliography}{99}

\bibitem{who2023}
World Health Organization, ``Asthma Fact Sheet,'' WHO Media Centre, 2023.
[Online]. Available: \url{https://www.who.int/news-room/fact-sheets/detail/asthma}

\bibitem{guarnieri2014}
M. Guarnieri and J. R. Balmes, ``Outdoor air pollution and asthma,''
\textit{The Lancet}, vol.~383, no.~9928, pp.~1581--1592, May 2014.

\bibitem{patel2019}
R. Patel, M. Smith, and J. Johnson, ``Machine learning approaches for predicting
asthma exacerbations from electronic health records,''
\textit{Journal of Asthma}, vol.~56, no.~8, pp.~842--851, 2019.

\bibitem{xiang2020}
Y. Xiang, L. Chen, and W. Zhang, ``Deep learning architectures for asthma risk
prediction,'' \textit{IEEE Access}, vol.~8, pp.~128565--128578, 2020.

\bibitem{finkelstein2017}
J. Finkelstein and I. Jeong, ``Machine learning approaches for early asthma risk
prediction,'' \textit{Annals of the New York Academy of Sciences},
vol.~1387, no.~1, pp.~153--165, Jan. 2017.

\bibitem{liu2019}
C. Liu, R. Chen, and S. Zhao, ``Ambient particulate air pollution and daily
mortality in major cities,'' \textit{New England Journal of Medicine},
vol.~381, pp.~705--715, Aug. 2019.

\bibitem{garcez2020}
A. S. d'Avila Garcez and L. C. Lamb, ``Neurosymbolic artificial intelligence:
The state of the art,'' \textit{arXiv preprint arXiv:2012.05876}, Dec. 2020.

\bibitem{caruana2015}
R. Caruana, Y. Lou, J. Gehrke, P. Koch, M. Sturm, and N. Elhadad, ``Intelligible
models for healthcare: Predicting pneumonia risk and hospital readmission,'' in
\textit{Proc. 21st ACM SIGKDD Int. Conf. Knowledge Discovery and Data Mining},
Sydney, Australia, 2015, pp.~1721--1730.

\bibitem{wolpert1992}
D. H. Wolpert, ``Stacked generalization,'' \textit{Neural Networks},
vol.~5, no.~2, pp.~241--259, 1992.

\bibitem{chen2016}
T. Chen and C. Guestrin, ``XGBoost: A scalable tree boosting system,'' in
\textit{Proc. 22nd ACM SIGKDD Int. Conf. Knowledge Discovery and Data Mining},
San Francisco, CA, USA, 2016, pp.~785--794.

\bibitem{ke2017}
G. Ke, Q. Meng, T. Finley, T. Wang, W. Chen, W. Ma, Q. Ye, and T.-Y. Liu,
``LightGBM: A highly efficient gradient boosting decision tree,'' in
\textit{Proc. 31st Int. Conf. Neural Information Processing Systems},
Long Beach, CA, USA, 2017, pp.~3149--3157.

\bibitem{haque2024}
R. Haque, S. Ahmed, and M. Rahman, ``Clinical Asthma Dataset for Machine Learning
Research,'' Zenodo, 2024. [Online]. Available: \url{https://zenodo.org/}

\bibitem{lundberg2017}
S. M. Lundberg and S.-I. Lee, ``A unified approach to interpreting model
predictions,'' in \textit{Proc. 31st Int. Conf. Neural Information Processing
Systems}, Long Beach, CA, USA, 2017, pp.~4768--4777.

\bibitem{breiman2001}
L. Breiman, ``Random forests,'' \textit{Machine Learning},
vol.~45, no.~1, pp.~5--32, Oct. 2001.

\bibitem{shickel2018}
B. Shickel, P. J. Tighe, A. Bihorac, and P. Rashidi, ``Deep EHR: A survey of
recent advances in deep learning techniques for electronic health record analysis,''
\textit{IEEE J. Biomed. Health Inform.}, vol.~22, no.~5, pp.~1589--1604,
Sep. 2018.

\bibitem{kaggleasthma2024}
``Asthma Disease Dataset,'' Kaggle, 2024. [Online]. Available:
\url{https://www.kaggle.com/datasets/rabieelkharoua/asthma-disease-dataset}

\end{thebibliography}

\end{document}
